\documentclass[11pt]{article}
\usepackage[spanish]{babel}
\usepackage[utf8]{inputenc}
\usepackage[T1]{fontenc}
\usepackage{amsmath,amssymb,booktabs}
\usepackage[margin=2.2cm]{geometry}
\usepackage{enumitem}
\setlength{\parindent}{0pt}

\begin{document}

\begin{center}
{\Large\bfseries Econometría II --- E3: P3 (a--b)}\\[2pt]
Scoring: WOE, IV y escalamiento \quad$\cdot$\quad 25 min con reloj \quad$\cdot$\quad dom 22-11
\end{center}

\hrule\medskip

Un banco tramifica el \textbf{número de consultas} de otras instituciones a la deuda del cliente en los últimos 6 meses, en una muestra de desarrollo de 10.000 créditos ($y=1$: \textbf{malo}).

\begin{center}
\renewcommand{\arraystretch}{1.15}
\begin{tabular}{lrr}\toprule
Consultas & Buenos & Malos\\\midrule
0 & 3.200 & 500\\
1 & 2.000 & 400\\
2-3 & 1.600 & 500\\
4-5 & 600 & 100\\
6 o más & 600 & 500\\\midrule
Total & 8.000 & 2.000\\\bottomrule
\end{tabular}
\end{center}

\begin{enumerate}[label=\textbf{\alph*)}, leftmargin=*, itemsep=10pt]
\item ¿Qué signo de relación espera entre consultas y riesgo, y por qué? Calcule la tasa de malos de la cartera y, para cada tramo, $\%B$, $\%M$, WOE $\big(\ln(\%B/\%M)\cdot 100\big)$ y tasa de malos. \hfill [7 pts]

\item Calcule el IV y clasifique el poder predictivo. \hfill [3 pts]
\end{enumerate}

\medskip\hrule\medskip

{\small
\textbf{Fórmulas.}\quad
$\%B_i = \dfrac{B_i}{8.000}$,\quad $\%M_i = \dfrac{M_i}{2.000}$,\quad
tasa de malos$_i = \dfrac{M_i}{B_i+M_i}$,\quad
$\text{IV} = \displaystyle\sum_i (\%B_i-\%M_i)\,\ln\!\frac{\%B_i}{\%M_i}$ \ (sin $\times 100$).

\smallskip
\textbf{Escala del IV.}\quad $<0{,}02$ inútil \ $\cdot$ \ $0{,}02$--$0{,}1$ débil \ $\cdot$ \ $0{,}1$--$0{,}3$ medio \ $\cdot$ \ $0{,}3$--$0{,}5$ fuerte \ $\cdot$ \ $>0{,}5$ sospechoso.
}

\bigskip
\textbf{Tabla de trabajo}

\begin{center}
\renewcommand{\arraystretch}{1.6}
\begin{tabular}{l|c|c|c|c|c|c}\toprule
Consultas & $\%B$ & $\%M$ & $\%B-\%M$ & WOE & aporte IV & tasa malos\\\midrule
0 & & & & & & \\
1 & & & & & & \\
2-3 & & & & & & \\
4-5 & & & & & & \\
6 o más & & & & & & \\\midrule
Total & 1 & 1 & 0 & --- & IV $=$ & \\\bottomrule
\end{tabular}
\end{center}

\end{document}
